\documentclass[submission,copyright,creativecommons]{eptcs}
\providecommand{\event}{ACT 2026} % Applied Category Theory 2026

\usepackage{iftex}

\ifpdf
  \usepackage{underscore}
  \usepackage[T1]{fontenc}
\fi

% Math packages
\usepackage{amsmath}
\usepackage{amssymb}
\usepackage{amsthm}

% Graphics
\usepackage{graphicx}

% Tables
\usepackage{booktabs}

% Code listings
\usepackage{listings}
\lstset{
  language=Haskell,
  basicstyle=\small\ttfamily,
  keywordstyle=\bfseries,
  columns=flexible,
  keepspaces=true
}

% Tikz for diagrams
\usepackage{tikz}
\usetikzlibrary{arrows.meta,positioning,cd}

% Title and authors
\title{From Games to Graphs:\\Categorical Composition of Genetic Algorithms Across Domains}

\author{Lyra Vega
\institute{Open University Helsinki}
\email{lyraclaude20@gmail.com}
\and
Claudius Turing
\institute{Independent Researcher}
\and
Robin Langer
\institute{Independent Researcher}
}

\newcommand{\titlerunning}{From Games to Graphs}
\newcommand{\authorrunning}{L. Vega, C. Turing \& R. Langer}

\hypersetup{
  bookmarksnumbered,
  pdftitle    = {\titlerunning},
  pdfauthor   = {\authorrunning},
  pdfsubject  = {ACT 2026},
  pdfkeywords = {category theory, genetic algorithms, Kleisli morphisms, evolutionary computation, compositional optimization}
}

% Theorem environments
\newtheorem{theorem}{Theorem}
\newtheorem{conjecture}{Conjecture}
\newtheorem{definition}{Definition}
\newtheorem{proposition}{Proposition}
\newtheorem{lemma}{Lemma}
\newtheorem{corollary}{Corollary}
\newtheorem{remark}{Remark}
\newtheorem{example}{Example}

\begin{document}
\maketitle

\begin{abstract}
Four independent research programs have recently converged on the same
structural insight: in optimization, \emph{how} components compose determines
system-level behavior more than the components themselves. Gavranovi\'{c} et
al.\ formalized neural network architectures categorically; Hedges and
Sakamoto cast reinforcement learning in categorical cybernetics; Bakirtzis et
al.\ derived safety guarantees for compositional MDPs. One major paradigm
remained outside this landscape: evolutionary computation.

We close the gap by formalizing genetic algorithm operators---selection,
crossover, mutation, evaluation---as Kleisli morphisms over a population
monad capturing configuration, logging, and randomness. Pipelines, strategies,
and island models arise as compositions at three hierarchical levels in the
Kleisli category, providing the first categorical formalization of genetic
algorithms. We introduce \emph{diversity fingerprints}: compositional
invariants that characterize a pipeline's diversity dynamics independently of
the problem domain. Experiments on two unrelated domains (checkers strategy
evolution and maze topology optimization) show that composition pattern, not
domain content, determines the fingerprint, with strict composition preserving
$2.2\times$ higher diversity than lax (Cohen's $d = 4.34$,
$p < 3.69 \times 10^{-11}$). We conjecture that strict composition preserves
compositional invariants while lax composition degrades them---a pattern
supported by five independent studies across optimization paradigms.
\end{abstract}

%% ============================================================
%% Section 1: The Optimization Zoo (~2 pages)
%% ============================================================
\section{Introduction: Four Groups, One Insight}
\label{sec:intro}

Over the past three years, four independent research groups---working on four
different optimization paradigms---have arrived at the same structural insight:
the way an optimization pipeline \emph{composes} its components determines its
dynamics more than the components themselves. This convergence is not
coincidental. It reflects a deep mathematical regularity that category theory
is uniquely positioned to articulate.

\paragraph{Neural networks.}
Gavranovi\'{c} et al.~\cite{gavranovic2024categorical} showed that deep learning
architectures can be understood as constructions in the category $\mathbf{Para}$,
where gradient descent is a functor and backpropagation emerges as its
composition law. The key insight is not that individual layers are interesting
categorically---it is that the \emph{way layers compose} (via the chain rule,
formalized as Kleisli composition for a suitable monad) determines both the
expressiveness and the trainability of the resulting network. Their framework
was presented at ICML~2024 and has since informed the design of compositional
neural architectures at scale.

\paragraph{Reinforcement learning.}
Hedges and Sakamoto~\cite{hedges2024reinforcement} formalized RL within
categorical cybernetics (EPTCS~429), expressing the Bellman equation as
a natural transformation. The compositional structure of the
agent--environment loop, not the choice of value function, governs
convergence.

\paragraph{Compositional RL.}
In JMLR~v26, Bakirtzis et al.~\cite{bakirtzis2024compositional}
formalized Markov decision processes categorically, deriving safety
guarantees from composition structure. Their approach is the closest
analogue to ours: formalize an optimization paradigm categorically
and show that composition reveals structure invisible at the
component level.

\paragraph{Evolutionary computation: the remaining gap.}
One major optimization paradigm has remained outside this categorical
landscape: evolutionary computation. Genetic algorithms have been studied
extensively for over five decades~\cite{holland1975adaptation,goldberg1989genetic,
dejong2006evolutionary}, yet their formalization has remained largely
set-theoretic or statistical. The compositional structure of a GA pipeline---selection,
crossover, mutation, composed in sequence---has been treated as an
implementation detail rather than a mathematical object worthy of study in its
own right. This paper addresses that gap.

\medskip

The common thread is that \emph{how} components compose predicts
system-level behavior invisible at the component level. This
observation has a natural categorical formulation. When composition is
\emph{strict}---meaning it respects component boundaries exactly, with
information flowing only through defined interfaces---compositional invariants
are preserved. When composition is \emph{lax}---meaning additional communication
occurs between components beyond the formal interface---those invariants may
degrade. This strict/lax distinction is well-studied in category theory (a
strict monoidal functor preserves structure on the nose; a lax monoidal functor
preserves it only up to coherent natural transformations), and it connects our
work to Ghani et al.'s recent formalization of approximate game
theory~\cite{ghani2025approximate}, where exact Nash equilibria compose cleanly
but approximate equilibria do not. We conjecture that this pattern is general:
compositional invariants are preserved under strict composition and disrupted
under lax composition, across optimization paradigms.

Our formalization proceeds via Moggi's framework of monads and Kleisli
categories~\cite{moggi1991notions}. We model the state of an evolutionary
computation as a monad $T = \mathit{Reader} \times \mathit{Writer} \times
\mathit{State}$, capturing configuration, logging, and randomness respectively.
GA operators become Kleisli morphisms $a \to T\,b$, and the composition of
operators into pipelines is Kleisli composition. This is not merely a
notational convenience. By expressing GA pipelines as composed Kleisli
morphisms, we gain access to the full machinery of categorical
reasoning---functors, natural transformations, and compositional invariants---and
can state and test precise claims about how pipeline structure determines
evolutionary dynamics.

\paragraph{Contributions.}
This paper makes three contributions:
\begin{enumerate}
\item We formalize genetic algorithm operators as Kleisli morphisms over a
  population monad, completing the categorical optimization landscape initiated
  by~\cite{gavranovic2024categorical,hedges2024reinforcement,bakirtzis2024compositional}.
\item We introduce \emph{diversity fingerprints}---compositional invariants that
  characterize the diversity dynamics of a GA pipeline---and demonstrate
  empirically that the same composition produces the same fingerprint across
  two unrelated problem domains (game-playing strategy evolution and maze
  topology optimization), with effect size $d = 4.34$ and $p < 3.69 \times
  10^{-11}$.
\item We propose a general conjecture: across optimization paradigms,
  compositional structure predicts system-level behavior. Specifically, strict
  composition preserves compositional invariants while lax composition
  disrupts them.
\end{enumerate}

\paragraph{Outline.}
Section~\ref{sec:bridge} describes the disconnected state of the evolutionary
computation literature and the bridging role of categorical vocabulary.
Section~\ref{sec:framework} presents our Kleisli formalization of GA operators.
Section~\ref{sec:domains} instantiates the framework on two domains.
Section~\ref{sec:fingerprints} presents the diversity fingerprint results.
Section~\ref{sec:conjecture} states our general conjecture with supporting
evidence from all four categorical optimization programs.
Section~\ref{sec:discussion} discusses limitations and future directions.

%% ============================================================
%% Section 2: The Knowledge Graph Hole (~1 page)
%% ============================================================
\section{Bridging Disconnected Communities}
\label{sec:bridge}

The convergence described in Section~\ref{sec:intro} is invisible from
within any single community. To understand why, we constructed a knowledge
graph of the evolutionary computation literature. Our corpus comprises 76
papers spanning five research communities---neuroevolution, morphology-control
co-evolution, open-ended evolution, quality-diversity, and game-theoretic
coevolution---yielding 555 concept nodes and 446 edges, where an edge
connects two concepts that co-occur in the same paper.

\paragraph{The disconnection.}
The graph reveals a striking structural hole. Game-theoretic coevolution
shares \emph{zero} concept overlap with three of the four other clusters.
Its only bridge to the rest of the corpus is the term ``fitness,'' shared
with morphology-control---a term so generic that it provides no structural
connection. This isolation is not an artifact of our corpus selection. It
reflects a genuine disciplinary fragmentation: researchers studying the
evolution of competitive game strategies and researchers evolving robot
morphologies use the same algorithmic machinery---population-based search
with selection, crossover, and mutation---yet cite different foundational
works, attend different conferences, and employ different vocabularies.

The disconnection is linguistic, not algorithmic. Consider migration:
in island-model genetic algorithms~\cite{whitley1999island,skolicki2005influence},
subpopulations exchange individuals periodically to maintain diversity.
In game-theoretic coevolution, the same mechanism appears as
``strategy transfer'' between competing populations. In quality-diversity
archives, it surfaces as ``niche injection.'' The underlying operation---a
functorial map transporting individuals between subpopulations---is
identical in all three cases. But because each community names it
differently, the shared structure remains invisible in the literature graph.

\paragraph{Categorical vocabulary as bridge.}
This paper introduces categorical vocabulary that makes the shared
structure explicit. In our framework (Section~\ref{sec:framework}), each
GA operator is a Kleisli morphism $a \to T\,b$ over a population monad $T$.
This single formalization captures what the communities already do under
different names:
\begin{itemize}
\item \textbf{Mutation} becomes an endomorphism $a \to T\,a$: a Kleisli
  arrow whose domain and codomain coincide.
\item \textbf{Selection} becomes a projection: a morphism that reduces the
  population by discarding individuals according to a fitness criterion.
\item \textbf{Crossover} becomes a product construction: a morphism that
  combines information from two parent individuals into offspring.
\item \textbf{Migration} becomes a functorial transformation: a natural
  transformation between population functors on different subpopulations.
\end{itemize}
These are not new algorithms. They are new \emph{names} for existing
algorithms---names drawn from a mathematical vocabulary (category theory)
that is designed precisely for describing how components compose. The
bridge we build is linguistic: we do not propose better selection operators
or faster crossover schemes. We provide a shared formal language in which
the structural identity across communities becomes a provable fact rather
than an informal analogy.

\paragraph{Plugging the hole.}
In the knowledge graph, bridge papers are those that span structural
holes between otherwise disconnected clusters. The most effective bridges
in our corpus are NEAT~\cite{stanley2002neat}
(bridging neuroevolution to multiple communities) and novelty
search~\cite{lehman2011novelty}
(bridging open-ended evolution to quality-diversity). Our contribution
bridges the deepest hole: \textbf{game-theoretic coevolution to
structural optimization}. The categorical vocabulary of Kleisli morphisms,
functors, and natural transformations provides exactly the shared
terminology that the knowledge graph lacks. By introducing this vocabulary,
we do not merely describe a connection that already existed informally---we
create new edges in the knowledge graph, connecting concepts that were
previously isolated despite denoting the same underlying operations.

%% ============================================================
%% Section 3: Framework (~2 pages)
%% ============================================================
\section{Genetic Algorithm Operators as Kleisli Morphisms}
\label{sec:framework}

We formalize the operators of a genetic algorithm as Kleisli morphisms
over a monad that captures the computational effects inherent in
evolution: configuration, logging, and randomness. The key ideas are
due to Moggi~\cite{moggi1991notions}, who showed that computational
effects---side-effects, nondeterminism, exceptions---are uniformly
captured by monads. Our contribution is the observation that
\emph{evolutionary dynamics are computational effects too}: the
stochastic perturbation of mutation, the stateful threading of a
pseudorandom generator, the accumulation of diversity metrics over
generations. The monad $T$ makes these effects explicit and compositional.

\subsection{The Evolution Monad}
\label{sec:monad}

\begin{definition}[Population]
Let $\Sigma$ be a type of \emph{genomes} (e.g., $\mathbb{R}^n$ for
weight vectors, $\{0,1\}^k$ for bitstrings). A \emph{scored individual}
is a pair $(g, f) \in \Sigma \times \mathbb{R}$, where $f$ is the
fitness of genome~$g$. A \emph{population} is a finite list
$P = [(g_1, f_1), \ldots, (g_n, f_n)]$ of scored individuals.
We write $\mathsf{Pop}(\Sigma)$ for the type of populations over~$\Sigma$.
\end{definition}

\begin{definition}[Evolution monad]
\label{def:evo-monad}
The \emph{evolution monad} is the composite monad
\[
  T \;=\; \mathsf{Reader}(\Gamma) \times \mathsf{Writer}(\Lambda) \times \mathsf{State}(S),
\]
where:
\begin{itemize}
\item $\Gamma$ is the \emph{configuration context}: population size,
  mutation rate, crossover rate, tournament size---read-only parameters
  accessed by every operator but never modified.
\item $\Lambda$ is a \emph{log monoid} of diversity metrics: each
  generation appends a record (best fitness, mean fitness, diversity
  index) to an accumulated history.  The monoid structure
  $(\Lambda, \cdot, \varepsilon)$ ensures that logs from composed
  operators concatenate associatively.
\item $S$ is the state of a pseudorandom number generator, threaded
  through every stochastic operation (selection, crossover, mutation).
\end{itemize}
Concretely, for a result type $A$, a computation in~$T$ is a function
\[
  \Gamma \to S \to (A \times S \times \Lambda),
\]
reading a configuration, consuming and producing PRNG state, and
emitting a log. The unit $\eta_A : A \to T\,A$ lifts a pure value
into~$T$ by returning the configuration and state unchanged with the
empty log~$\varepsilon$. The multiplication $\mu_A : T^2\,A \to T\,A$
sequences two $T$-computations by threading the PRNG state and
concatenating the logs.
\end{definition}

\subsection{Operators as Kleisli Morphisms}
\label{sec:kleisli-ops}

\begin{definition}[Genetic operator]
A \emph{genetic operator} is a Kleisli morphism for~$T$:
an arrow $f : \mathsf{Pop}(A) \to T(\mathsf{Pop}(B))$, where $A$~and~$B$
are genome types (possibly equal). We write $\mathbf{Kl}(T)$ for the
Kleisli category whose objects are population types and whose morphisms
are genetic operators.
\end{definition}

The standard GA operators are Kleisli morphisms:
\begin{itemize}
\item \textbf{Evaluation.} $\mathsf{eval}_\phi : \mathsf{Pop}(\Sigma) \to
  T(\mathsf{Pop}(\Sigma))$, parameterized by a fitness function
  $\phi : \Sigma \to \mathbb{R}$. Pairs each genome with its fitness.

\item \textbf{Selection.} $\mathsf{sel} : \mathsf{Pop}(\Sigma) \to
  T(\mathsf{Pop}(\Sigma))$---an endomorphism projecting the
  population onto a fitter subset (consuming PRNG state, logging
  diversity).

\item \textbf{Crossover.} $\mathsf{cross} : \mathsf{Pop}(\Sigma) \to
  T(\mathsf{Pop}(\Sigma))$---recombines genetic material from pairs
  of individuals at a random point, with probability read from~$\Gamma$.

\item \textbf{Mutation.} $\mathsf{mut}_\delta : \mathsf{Pop}(\Sigma)
  \to T(\mathsf{Pop}(\Sigma))$---an endomorphism parameterized by
  $\delta : \Sigma \to T(\Sigma)$
  (Gaussian noise for $\mathbb{R}^n$, bit-flip for $\{0,1\}^k$).

\item \textbf{Migration.} A natural transformation
  $\prod_{i \in I} \mathsf{Pop}(\Sigma) \to
  \prod_{i \in I} T(\mathsf{Pop}(\Sigma))$
  transporting individuals between subpopulations without inspecting
  them. Its naturality---commuting with any operator applied per
  island---underlies the strict/lax dichotomy (Section~\ref{sec:conjecture}).
\end{itemize}

\begin{definition}[Kleisli composition]
\label{def:kleisli-comp}
Given $f : A \to T\,B$ and $g : B \to T\,C$,
their Kleisli composite $g \circ_K f : A \to T\,C$ is
$g \circ_K f = \mu_C \circ T(g) \circ f$,
where $\mu : T^2 \Rightarrow T$ is the monad multiplication.
The monad laws guarantee associativity and identity, so
$\mathbf{Kl}(T)$ is a category.
\end{definition}

\subsection{Three-Level Composition}
\label{sec:three-level}

The Kleisli category $\mathbf{Kl}(T)$ supports composition at three
distinct levels:

\paragraph{Level 1: Operators $\to$ Pipelines.}
A \emph{pipeline} (one generation) is the Kleisli composite of
individual operators:
\[
  \pi \;=\; \mathsf{mut}_\delta \circ_K \mathsf{cross}
    \circ_K \mathsf{sel} \circ_K \mathsf{eval}_\phi.
\]
The pipeline $\pi : \mathsf{Pop}(\Sigma) \to T(\mathsf{Pop}(\Sigma))$
is itself a Kleisli morphism---a single arrow encoding one complete
generation. Running $n$~generations is the $n$-fold self-composition
$\pi^n = \pi \circ_K \pi \circ_K \cdots \circ_K \pi$.

\paragraph{Level 2: Pipelines $\to$ Strategies.}
A \emph{strategy} composes pipelines with different parameters to
implement multi-phase evolution. For example, the \emph{hourglass}
strategy composes an exploration phase (high mutation, large
tournament) with a convergence phase (low mutation, small tournament)
and a re-diversification phase:
\[
  \sigma_{\text{hg}} \;=\; \pi_{\text{diversify}}^{25}
    \circ_K \pi_{\text{converge}}^{10}
    \circ_K \pi_{\text{explore}}^{15}.
\]
The sequential composition of strategies is again associative.
Conditional composition (\emph{adaptive} strategies) introduces a
predicate on intermediate results: run~$\sigma_1$, test a condition
(e.g., plateau detection), and either accept the result or continue
with~$\sigma_2$.

\paragraph{Level 3: Strategies $\to$ Multi-population systems.}
The \emph{island functor} $\mathcal{I}$ lifts a strategy on a single
population to a strategy on $n$~independent subpopulations with
periodic migration. Formally, $\mathcal{I}$ maps objects
$\mathsf{Pop}(\Sigma) \mapsto \mathsf{Pop}(\Sigma)^n$ and morphisms
$\sigma \mapsto \mathcal{I}(\sigma)$, where $\mathcal{I}(\sigma)$
applies~$\sigma$ to each island independently and interleaves
migration events. Whether $\mathcal{I}$ preserves sequential
composition---i.e., whether $\mathcal{I}(\sigma_1 \circ_K \sigma_2) =
\mathcal{I}(\sigma_1) \circ_K \mathcal{I}(\sigma_2)$---depends on
whether any migration occurs (Section~\ref{sec:conjecture}).

\medskip
At each level, the composite has the same type as its parts: a
pipeline of operators is itself an operator; a sequence of strategies
is itself a strategy; an island model containing independent GAs is
itself a GA. Composition closes over the same categorical structure.

\subsection{Implementation}
\label{sec:haskell}

The categorical structure is directly implemented in Haskell.
The evolution monad, genetic operator type, and Kleisli composition are:

\begin{lstlisting}
type EvoM = ReaderT GAConfig
            (StateT StdGen (Writer GALog))

newtype GeneticOp m a b =
  GeneticOp { runOp :: [a] -> m [b] }

(>>>:) :: Monad m => GeneticOp m a b
       -> GeneticOp m b c -> GeneticOp m a c
(GeneticOp f) >>>: (GeneticOp g) =
  GeneticOp (\pop -> f pop >>= g)
\end{lstlisting}

\noindent
A complete pipeline composes operators left-to-right:

\begin{lstlisting}
pipeline :: GeneticOp EvoM [a] [a]
pipeline = evaluate fitnessFunc
  >>>: elitistSelect
  >>>: onePointCrossover
  >>>: pointMutate mutFunc
\end{lstlisting}

\noindent
Because \texttt{pipeline} is itself a \texttt{GeneticOp}, it plugs
into any strategy combinator---\texttt{sequential}, \texttt{race},
\texttt{adaptive}, \texttt{islandStrategy}---without modification.
Changing the fitness function or mutation operator changes the
\emph{content}; the compositional \emph{structure} is invariant.
The same strategy code operates unchanged on checkers weight vectors
(\texttt{[Double]}), maze topologies (\texttt{[Bool]}), and expression
trees (\texttt{ExprTree}).

%% ============================================================
%% Section 4: Two Domains (~2 pages)
%% ============================================================
\section{Two Domains, One Category}
\label{sec:domains}

We instantiate the framework on two domains chosen for maximum
dissimilarity: checkers (adversarial game strategy, real-valued
genomes) and mazes (structural design, binary genomes). The point is
not that either domain is novel, but that the categorical pipeline is
\emph{identical} across both. Before describing the domains, we define
the selection mechanism used in all our experiments, since it
illustrates how a concrete algorithmic operation becomes a Kleisli
morphism.

\begin{definition}[Tournament selection]
\label{def:tournament}
Let $P = [(g_1,f_1),\ldots,(g_n,f_n)]$ be a scored population and
$k \in \mathbb{N}$ the \emph{tournament size} (read from the
configuration context~$\Gamma$). \emph{Tournament selection} draws
$k$~individuals uniformly at random from~$P$ and returns the one with
the highest fitness. This is repeated to fill a new population of a
desired size.

Formally, $\mathsf{sel}_k : \mathsf{Pop}(\Sigma) \to
T(\mathsf{Pop}(\Sigma))$ is a Kleisli endomorphism that reads $k$
from~$\Gamma$, consumes PRNG state from~$S$ to sample the
$k$-subsets, and logs the diversity of the selected population
to~$\Lambda$.
\end{definition}

Larger $k$ means stronger selection pressure ($k = n$ always selects
the best individual; $k = 1$ selects uniformly at random). With
tournament selection defined, we describe the two domains.

\subsection{Checkers: Evolving Evaluation Functions}
\label{sec:checkers}

Following Samuel's pioneering work~\cite{samuel1959checkers}, we
evolve a population of 30 weight vectors, each encoding 10 features
of an 8$\times$8 checkers board position: material advantage, king
count, back-row defense, center control, advancement toward promotion,
mobility, vulnerability (pieces under threat), protection (pieces
defended by allies), king centrality, and edge positioning. Each
weight $w_i \in [-5, 5]$ determines the relative importance of
feature~$i$ in a linear evaluation function $\phi(\mathbf{w}, b) =
\sum_i w_i \cdot x_i(b)$ for board position~$b$.

Fitness is win rate against two fixed opponent strategies (two games
each, alternating colors). Mutation is Gaussian perturbation
$\delta(\mathbf{w}) = \mathbf{w} + \mathcal{N}(0, \sigma^2 I)$;
crossover is one-point on the weight vector.

In the categorical framework, a single generation is the Kleisli
composite (Definition~\ref{def:kleisli-comp}):
\[
  \pi_{\mathrm{ck}} \;=\; \mathsf{mut}_\delta \circ_K \mathsf{cross}
    \circ_K \mathsf{sel}_k \circ_K \mathsf{eval}_\phi
  \;:\; \mathsf{Pop}(\mathbb{R}^{10}) \to T(\mathsf{Pop}(\mathbb{R}^{10})).
\]

\subsection{Mazes: Evolving Topological Structures}
\label{sec:mazes}

The maze domain evolves $6 \times 6$ grid mazes. The genome is a
60-bit binary vector encoding wall presence on each potential edge of
the grid graph. Fitness is a composite of three BFS-derived metrics:
solution path length (longer paths increase difficulty), dead-end
ratio (more dead ends create confusion), and branching factor (more
choice points make navigation interesting).

Mutation is bit-flip: $\delta(g)_i = 1 - g_i$ with probability $p$
per bit, read from~$\Gamma$. Crossover is one-point on the
binary vector.

The generation pipeline is:
\[
  \pi_{\mathrm{mz}} \;=\; \mathsf{mut}_\delta \circ_K \mathsf{cross}
    \circ_K \mathsf{sel}_k \circ_K \mathsf{eval}_\psi
  \;:\; \mathsf{Pop}(\{0,1\}^{60}) \to T(\mathsf{Pop}(\{0,1\}^{60})).
\]

The fitness function $\psi$ measures topological complexity; $\phi$
measures playing skill. The genome types are incompatible:
$\mathbb{R}^{10}$ versus $\{0,1\}^{60}$. The mutation operators are
domain-specific (Gaussian perturbation versus bit-flip). And yet the
\emph{pipeline shape} is identical: the same four Kleisli morphisms
composed in the same order.

\subsection{Four Strategies on Each Domain}
\label{sec:four-strategies}

We run four composition strategies on each domain, each corresponding
to a distinct way of composing pipelines in~$\mathbf{Kl}(T)$:

\begin{enumerate}
\item \textbf{Flat generational.} The pipeline~$\pi$ applied for
  50~generations: $\pi^{50}$. A single morphism iterated.

\item \textbf{Hourglass.} A three-phase sequential composition:
  $\sigma_{\mathrm{hg}} = \pi_{\mathrm{div}}^{25} \circ_K
  \pi_{\mathrm{conv}}^{10} \circ_K \pi_{\mathrm{exp}}^{15}$, where
  the phases use different mutation rates and tournament sizes
  (exploration: high mutation, large $k$; convergence: low mutation,
  small $k$; diversification: medium mutation, large $k$).

\item \textbf{Island model.} Four subpopulations with ring-topology
  migration every 5~generations, implemented via the island
  functor~$\mathcal{I}$ from Section~\ref{sec:three-level}.

\item \textbf{Adaptive.} Exploration with plateau detection: run
  $\pi_{\mathrm{exp}}$, monitor fitness improvement, and switch to
  $\pi_{\mathrm{focus}}$ when improvement stalls for 5 consecutive
  generations. This is a conditional composition
  $\sigma_{\mathrm{ad}} = \mathsf{cond}(\mathsf{plateau},
  \pi_{\mathrm{focus}}, \pi_{\mathrm{exp}})$.
\end{enumerate}

\subsection{What Changes, What Stays}
\label{sec:invariance}

\begin{table}[t]
\centering
\caption{Invariance across domains. The categorical structure (pipeline
shape, monad stack, composition operator, strategy combinators) is
identical; only the domain-specific content (genome type, fitness
function, mutation operator) changes.}
\label{tab:invariance}
\begin{tabular}{@{}lllc@{}}
\toprule
Aspect & Checkers & Mazes & Invariant? \\
\midrule
Genome type & $\mathbb{R}^{10}$ & $\{0,1\}^{60}$ & No \\
Fitness function & Win rate & BFS composite & No \\
Mutation & Gaussian & Bit-flip & No \\
\midrule
Pipeline shape & $\mathsf{mut} \circ_K \mathsf{cross} \circ_K
  \mathsf{sel} \circ_K \mathsf{eval}$ & Same & \textbf{Yes} \\
Monad stack & $T = R \times W \times S$ & Same & \textbf{Yes} \\
Strategy combinators & hourglass, island, adaptive & Same & \textbf{Yes} \\
Composition operator & $\circ_K$ & Same & \textbf{Yes} \\
Category laws & Assoc., identity & Same & \textbf{Yes} \\
\bottomrule
\end{tabular}
\end{table}

Table~\ref{tab:invariance} summarizes the invariance. This is exactly
what category theory predicts: the structure lives in the arrows and
their composition, not in the objects. Changing the objects (genome
types) is a change of \emph{content}; the categorical \emph{structure}
is preserved. The question then becomes: does this structural
invariance produce measurable, stable consequences in the evolutionary
dynamics? The answer is yes, and the evidence is the subject of the
next section.

%% ============================================================
%% Section 5: Key Result (~2 pages)
%% ============================================================
\section{Diversity Fingerprints: Composition Determines Dynamics}
\label{sec:fingerprints}

We now present the central empirical result of this paper. Running
the four strategies from Section~\ref{sec:four-strategies} on both
domains produces \emph{diversity fingerprints}---characteristic
diversity trajectories that are determined by the composition pattern
and stable across unrelated problem domains.

\subsection{Fingerprints Defined}
\label{sec:fp-def}

\begin{definition}[Diversity fingerprint]
\label{def:fingerprint}
Let $\sigma : \mathsf{Pop}(\Sigma) \to T(\mathsf{Pop}(\Sigma))$ be a
strategy (a composed Kleisli morphism). The \emph{diversity fingerprint}
of~$\sigma$ is the sequence
$\mathbf{d}(\sigma) = (d_0, d_1, \ldots, d_n) \in \mathbb{R}^{n+1}$,
where $d_i$ is the genotypic diversity of the population at
generation~$i$. Genotypic diversity is measured as pairwise Hamming
distance (binary genomes) or pairwise Euclidean distance (real-valued
genomes), normalized to $[0, 1]$.
\end{definition}

A fingerprint captures the \emph{shape} of the diversity trajectory:
whether it declines monotonically, spikes and crashes, oscillates
steadily, or collapses. The claim is that this shape is determined by
the composition pattern~$\sigma$---how the Kleisli morphisms are
assembled---not by the genome type~$\Sigma$ or the fitness function.

\subsection{Four Fingerprints Observed}
\label{sec:fp-observed}

Table~\ref{tab:fingerprints} shows genotypic diversity over time for
each strategy on both domains. Four qualitatively distinct fingerprints
emerge:

\begin{table}[t]
\centering
\caption{Genotypic diversity (normalized pairwise distance) over time
for three of four strategies tested on both domains. Diversity is
sampled every 5~generations. Three distinct trajectory shapes emerge,
each consistent across domains. The fourth strategy (adaptive) was
tested only on symbolic regression; its spike--collapse fingerprint is
discussed in the text but omitted from this table.}
\label{tab:fingerprints}
\begin{tabular}{@{}r rr rr rr@{}}
\toprule
& \multicolumn{2}{c}{Flat} & \multicolumn{2}{c}{Hourglass}
& \multicolumn{2}{c}{Island} \\
\cmidrule(lr){2-3} \cmidrule(lr){4-5} \cmidrule(lr){6-7}
Gen & Ck & Mz & Ck & Mz & Ck & Mz \\
\midrule
0  & 0.49 & 0.47 & 0.49 & 0.47 & 0.49 & 0.47 \\
5  & 0.34 & 0.21 & 0.32 & 0.51 & 0.23 & 0.37 \\
10 & 0.33 & 0.26 & 0.33 & 0.29 & 0.47 & 0.31 \\
15 & 0.31 & 0.26 & 0.32 & 0.46 & 0.44 & 0.28 \\
20 & 0.30 & 0.27 & 0.30 & 0.46 & 0.34 & 0.32 \\
\bottomrule
\end{tabular}
\end{table}

\begin{enumerate}
\item \textbf{Flat: monotonic decline.} Diversity decreases steadily
  toward a low equilibrium. No structural intervention prevents
  convergence. Both domains show the same smooth decline:
  $0.49 \to 0.30$ (checkers), $0.47 \to 0.27$ (mazes).

\item \textbf{Hourglass: spike--crash--rebound.} The exploration phase
  drives diversity up, the convergence phase crashes it, and the
  diversification phase partially restores it. The three-phase structure
  of the Kleisli composite is directly visible in the trajectory. In
  mazes, the rebound is particularly clear: $0.47 \to 0.51 \to 0.29
  \to 0.46$.

\item \textbf{Island: stable oscillation.} After an initial
  adjustment, diversity settles into a band maintained by
  migration--isolation dynamics. The island functor
  (Section~\ref{sec:three-level}) maintains higher equilibrium
  diversity than the flat strategy in both domains.

\item \textbf{Adaptive: spike--collapse.} On a third domain
  (symbolic regression with expression trees), the adaptive strategy
  drives diversity from 22.98 to 52.88 during exploration, then
  collapses to 1.92 as plateau detection triggers premature convergence.
  No recovery mechanism exists.  This fourth fingerprint was not tested
  on checkers and mazes, but the pattern---conditional composition
  producing irreversible diversity loss---is structurally distinct from
  the other three.
\end{enumerate}

\subsection{Cross-Domain Stability}
\label{sec:fp-stability}

The key observation is that each fingerprint \emph{shape} is stable
across domains, despite the domains differing in every
content-level detail. The flat strategy always declines monotonically.
The hourglass always shows spike--crash--rebound. The island always
maintains a stable band. The adaptive always collapses. These
qualitative patterns hold across real-valued weight vectors
($\mathbb{R}^{10}$) and binary maze genomes ($\{0,1\}^{60}$), across
adversarial fitness (win rate) and intrinsic fitness (topological
complexity), across Gaussian mutation and bit-flip mutation.

This stability has a categorical explanation. The diversity
fingerprint is a function of the \emph{composition pattern}---the
sequence and nesting of Kleisli morphisms---and the composition
pattern is exactly what is invariant across domains
(Table~\ref{tab:invariance}). In categorical terms, the fingerprint
is \emph{functorial}: it respects composition.

\begin{proposition}[Fingerprint functoriality]
\label{prop:functorial}
Let $F : \mathbf{Kl}(T_{\mathrm{ck}}) \to \mathbf{Kl}(T_{\mathrm{mz}})$
be the domain-change functor that replaces the genome type, fitness
function, and mutation operator while preserving the composition
structure. Then for any strategy~$\sigma$ in the checkers domain,
the qualitative shape of $\mathbf{d}(F(\sigma))$ matches the shape
of $\mathbf{d}(\sigma)$. In particular, the ordering of final
diversities across the three cross-domain strategies is preserved:
$\mathbf{d}_n(\sigma_{\mathrm{flat}})
< \mathbf{d}_n(\sigma_{\mathrm{isl}}) <
\mathbf{d}_n(\sigma_{\mathrm{hg}})$.
\end{proposition}

\begin{proof}[Proof sketch]
$F$ preserves composition structure by construction, acting as the
identity on morphisms while replacing objects and operator content.
Since fingerprints are traces of the log component $\Lambda$
accumulated via Kleisli composition
(Definition~\ref{def:evo-monad}), and $F$ preserves the monad
multiplication $\mu$, the fingerprint shape is determined by the
composition pattern alone. The diversity ordering then follows from
the monotonic relationship between composition complexity (number of
diversity-restoring phases) and final diversity.
\end{proof}

This ordering---flat $<$ island $<$ hourglass for final genotypic
diversity---holds in both domains (Table~\ref{tab:fingerprints}).
The composition pattern, not the content, determines the diversity
dynamics.

\subsection{Statistical Confirmation}
\label{sec:fp-stats}

To move from qualitative observation to statistical rigor, we ran
30~independent seeds comparing strict composition (no migration,
four independent subpopulations) against lax composition (migration
every 5~generations, same four subpopulations) on 20-bit
OneMax~\cite{mitchell1992royal} for 100~generations.

\begin{definition}[Cohen's $d$]
Cohen's $d$ measures effect size as the difference in means divided by
the pooled standard deviation:
$d = (\bar{x}_1 - \bar{x}_2) / s_p$.
By convention, $d > 0.8$ is a ``large'' effect.
\end{definition}

\begin{table}[t]
\centering
\caption{Strict vs.\ lax composition over 30~seeds. Diversity is
  normalized pairwise Hamming distance.  By generation~50, the
  distributions barely overlap ($d = 2.55$). By generation~100,
  they are completely separated ($d = 4.34$, $p = 3.7 \times
  10^{-11}$).  Both conditions converge to the same optimal fitness.}
\label{tab:stats}
\begin{tabular}{@{}r cc cc@{}}
\toprule
Gen & Strict (mean $\pm$ std) & Lax (mean $\pm$ std) & $d$ & $p$ \\
\midrule
0   & $0.500 \pm 0.001$ & $0.500 \pm 0.001$ & 0.00 & 1.00 \\
25  & $0.117 \pm 0.014$ & $0.099 \pm 0.019$ & 1.05 & $4.0 \times 10^{-4}$ \\
50  & $0.048 \pm 0.008$ & $0.028 \pm 0.008$ & 2.55 & $4.2 \times 10^{-10}$ \\
75  & $0.037 \pm 0.004$ & $0.017 \pm 0.004$ & 5.36 & $3.0 \times 10^{-11}$ \\
100 & $0.035 \pm 0.005$ & $0.016 \pm 0.004$ & \textbf{4.34} & $\mathbf{3.7 \times 10^{-11}}$ \\
\bottomrule
\end{tabular}
\end{table}

Table~\ref{tab:stats} shows the results. By generation~25, the
separation is already significant ($d = 1.05$, $p < 5 \times
10^{-4}$). By generation~100, the effect size reaches $d = 4.34$ with
$p = 3.7 \times 10^{-11}$---the diversity distributions under strict
and lax composition do not overlap at all across 30~seeds. This is not
merely a ``large'' effect in the conventional sense; it is a
\emph{massive} effect, the kind rarely observed in empirical science.

Crucially, both conditions converge to the same optimal fitness
(20.0/20.0). The divergence is entirely in the diversity dynamics,
not in solution quality. Strict and lax compositions find equally good
answers via completely different population structures.

\subsection{Why This Result Is New}
\label{sec:fp-new}

Previous work on population diversity~\cite{burke2004diversity,crepinsek2013exploration}
measures diversity as a function of operator parameters (mutation rate,
selection pressure). Our result is different: the \emph{composition
pattern}---how operators are assembled into strategies---determines
diversity dynamics independently of the operators and the problem
domain. No prior work has demonstrated that GA composition patterns
produce stable, measurable invariants across unrelated domains. The
diversity fingerprint is a property of the Kleisli composite---a
categorical invariant of the pipeline structure.

%% ============================================================
%% Section 6: Conjecture (~0.5 pages)
%% ============================================================
\section{Towards a General Conjecture}
\label{sec:conjecture}

The fingerprint results of Section~\ref{sec:fingerprints} suggest a
pattern extending beyond genetic algorithms.

\begin{conjecture}[Strict composition preserves invariants]
\label{conj:strict-lax}
Let $\mathcal{C}$ be a monoidal category of optimization components.
Strict composition in $\mathcal{C}$---preserving component
boundaries exactly---preserves compositional invariants. Lax
composition---permitting information flow beyond formal
interfaces---degrades them.
\end{conjecture}

Five independent observations support this conjecture:
(1)~\emph{Our experiments} (Section~\ref{sec:fp-stats}): strict
composition preserves $2.2\times$ higher diversity than lax, with
$d = 4.34$~\cite{cohen1988statistical}.
(2)~\emph{Scaling laws}~\cite{chen2025scaling}: strict orchestration
of LLM agents shows $17.2\times$ error amplification versus
$4.4\times$ under lax composition.
(3)~\emph{Constitutional evolution}~\cite{kumar2026constitutional}:
near-strict agents ($0.9\%$ communication) outperform lax ($62.2\%$)
on alignment.
(4)~\emph{Evolved RL}~\cite{alphaevolve2026marl}: modular
(strict-like) agents generalize; monolithic (lax) agents overfit.
(5)~\emph{Semantic collapse}~\cite{alpay2025semantic}: unconstrained
message passing collapses diversity in emergent languages.

This pattern connects directly to Ghani et al.'s approximate game
theory~\cite{ghani2025approximate}, where exact equilibria compose
cleanly but $\varepsilon$-approximate equilibria do not---the same
lax degradation in a game-theoretic setting. A formal proof unifying
both instances is future work.

%% ============================================================
%% Section 7: Discussion (~0.3 pages)
%% ============================================================
\section{Discussion}
\label{sec:discussion}

The Kleisli formalization gives practitioners a Rosetta stone:
populations become $T$-algebras, operators become Kleisli morphisms,
pipelines become Kleisli composites, island models become functors,
and migration becomes a natural transformation. Diversity
fingerprints are the key practical tool: given a composition pattern,
one can predict qualitative diversity dynamics \emph{before} running
experiments on a new domain.

\paragraph{Limitations.}
We tested two domains; more are needed (sorting networks, symbolic
regression). The Haskell implementation is faithful but not
mechanically verified. Only Levels~1 and~3 of the composition
hierarchy (Section~\ref{sec:three-level}) are empirically
demonstrated. Conjecture~\ref{conj:strict-lax} has strong empirical
support but lacks a formal proof.

%% ============================================================
%% Section 8: Conclusion (~0.25 pages)
%% ============================================================
\section{Conclusion}
\label{sec:conclusion}

Four research groups---working on neural networks, reinforcement
learning, compositional MDPs, and now evolutionary computation---have
independently discovered that the way optimization components are
composed determines system-level behavior. We formalized the last
of these four paradigms: genetic algorithm operators as Kleisli
morphisms over a population monad, with diversity fingerprints as
the key compositional invariant. Our experiments show that the same
composition pattern produces the same fingerprint across unrelated
domains ($d = 4.34$, $p < 4 \times 10^{-11}$), and we conjecture
that the strict/lax distinction governing this phenomenon is
general across the optimization landscape.

Future work includes extending fingerprint analysis to additional
domains (sorting networks, program synthesis), pursuing a formal
proof of Conjecture~\ref{conj:strict-lax} via spectral methods on
the Kleisli composition operator, and exploring heterogeneous island
models where different subpopulations run different strategies. The
optimization zoo is converging. Category theory provides the
language in which that convergence becomes precise.

%% ============================================================
%% Acknowledgements
%% ============================================================
% \subparagraph*{Acknowledgements.}
% TODO: Add acknowledgements

%% ============================================================
%% Bibliography
%% ============================================================
\bibliographystyle{eptcs}
\bibliography{references}

\end{document}
